% ============================================================
%  Investigación: Historias de Usuario y Puntos de Función
%  Universidad de las Fuerzas Armadas "ESPE"
%
%  INSTRUCCIONES OVERLEAF:
%  1. Sube también "referencias.bib" al proyecto (mismo nivel que main.tex).
%  2. Menu → Compiler → selecciona "LuaLaTeX".
%  3. Haz clic en Recompile (puede requerir 2 compilaciones).
% ============================================================

\documentclass[12pt, a4paper]{article}

% ── Paquetes ─────────────────────────────────────────────────
\usepackage[utf8]{inputenc}
\usepackage[T1]{fontenc}
\usepackage[spanish]{babel}
\usepackage{geometry}
\usepackage{graphicx}
\usepackage{booktabs}
\usepackage{tabularx}
\usepackage{longtable}
\usepackage{array}
\usepackage{xcolor}
\usepackage{colortbl}
\usepackage{setspace}
\usepackage{parskip}
\usepackage{titlesec}
\usepackage{fancyhdr}
\usepackage{lmodern}
\usepackage{enumitem}
\usepackage{tcolorbox}
\usepackage{amsmath}
\usepackage{csquotes}
\usepackage{multirow}

% ── BibLaTeX con estilo APA ───────────────────────────────────
\usepackage[
  style    = apa,
  backend  = biber,
  language = spanish
]{biblatex}
\addbibresource{referencias.bib}

% ── Márgenes ─────────────────────────────────────────────────
\geometry{top=2.5cm, bottom=2.5cm, left=3cm, right=2.5cm}

% ── Colores ──────────────────────────────────────────────────
\definecolor{espeazul}{RGB}{0, 51, 102}
\definecolor{grisclaro}{RGB}{245, 245, 245}
\definecolor{grisoscuro}{RGB}{80, 80, 80}
\definecolor{azulclaro}{RGB}{220, 235, 255}
\definecolor{verdeclaro}{RGB}{220, 255, 220}
\definecolor{amarilloclaro}{RGB}{255, 250, 220}

% ── Estilo de secciones ───────────────────────────────────────
\titleformat{\section}{\large\bfseries\color{espeazul}}{\thesection.}{0.5em}{}[\titlerule]
\titleformat{\subsection}{\normalsize\bfseries\color{grisoscuro}}{\thesubsection.}{0.5em}{}
\titleformat{\subsubsection}{\small\bfseries\color{espeazul}}{\thesubsubsection.}{0.5em}{}

% ── Encabezado y pie ─────────────────────────────────────────
\pagestyle{fancy}
\fancyhf{}
\fancyhead[L]{\small\color{grisoscuro} Historias de Usuario y Puntos de Función}
\fancyhead[R]{\small\color{grisoscuro} ESPE · NRC 30744}
\fancyfoot[C]{\small\thepage}
\renewcommand{\headrulewidth}{0.4pt}

% ── Cajas ────────────────────────────────────────────────────
\tcbuselibrary{skins, breakable}
\newtcolorbox{ejemplobox}{colback=azulclaro, colframe=espeazul,
  fonttitle=\bfseries\small, title=Ejemplo, breakable}
\newtcolorbox{formulabox}{colback=verdeclaro, colframe=espeazul!60!green,
  fonttitle=\bfseries\small, title=Fórmula, breakable}
\newtcolorbox{notabox}{colback=amarilloclaro, colframe=orange!70!black,
  fonttitle=\bfseries\small, title=Nota importante, breakable}

\rowcolors{2}{grisclaro}{white}

% =============================================================
\begin{document}

% ── CARÁTULA ─────────────────────────────────────────────────
\begin{titlepage}
  \centering
  \vspace*{1cm}
  \includegraphics[width=5cm]{img/logoESPE}
  \vspace{1.2cm}

  {\Large\bfseries\color{espeazul} UNIVERSIDAD DE LAS FUERZAS ARMADAS ``ESPE''}
  \vspace{0.4cm}

  {\large\bfseries\color{grisoscuro} DEPARTAMENTO DE CIENCIAS DE LA COMPUTACIÓN}
  \vspace{2cm}

  \rule{0.85\textwidth}{0.6pt}\vspace{0.5cm}

  {\LARGE\bfseries\color{espeazul}
    Historias de Usuario\\[0.4em]
    \large y Puntos de Función (IFPUG)
  }
  \vspace{0.3cm}
  {\normalsize\color{grisoscuro} Investigación Académica}
  \vspace{0.5cm}

  \rule{0.85\textwidth}{0.6pt}\vspace{2cm}

  \begin{tabular}{rl}
    \textbf{Nombre:}  & Luis Cueva                            \\[6pt]
    \textbf{NRC:}     & 30744                                 \\[6pt]
    \textbf{Materia:} & Construcción y Evolución del Software \\[6pt]
    \textbf{Fecha:}   & 4 de mayo de 2026                     \\
  \end{tabular}

  \vfill
  {\small\color{grisoscuro} Sangolquí, Ecuador}
\end{titlepage}

% ── ÍNDICE ───────────────────────────────────────────────────
\newpage
\tableofcontents
\newpage
\setcounter{page}{1}
\onehalfspacing

% =============================================================
\section{Objetivos}

\subsection{Objetivo General}

Analizar y documentar las técnicas de Historias de Usuario y Puntos de Función
(IFPUG) como herramientas complementarias para la especificación de requisitos
y la estimación del tamaño funcional en proyectos de software ágiles, aplicadas
al sistema web de Gestión de Notas.

\subsection{Objetivos Específicos}

\begin{enumerate}[leftmargin=1.8cm, itemsep=3pt]
  \item Describir el origen, características y componentes de las Historias de
        Usuario dentro del marco de las metodologías ágiles.
  \item Explicar el método IFPUG de Puntos de Función, sus tipos de componentes
        y su proceso de cálculo paso a paso.
  \item Identificar la relación directa entre los elementos de una Historia de
        Usuario y los componentes funcionales del método IFPUG.
  \item Aplicar ambas técnicas al sistema de Gestión de Notas como caso de
        estudio práctico.
  \item Evaluar las ventajas, limitaciones y errores comunes en el uso de cada
        técnica.
\end{enumerate}

% =============================================================
\section{Marco Teórico: Metodologías Ágiles}

\subsection{El Manifiesto Ágil}

En febrero de 2001, un grupo de 17 desarrolladores firmaron el
\textbf{Manifiesto para el Desarrollo Ágil de Software}
\parencite{beck2001}, estableciendo cuatro valores fundamentales:

\vspace{0.3cm}
\begin{tabular}{p{6cm} p{6cm}}
  \toprule
  \rowcolor{espeazul}
  \textcolor{white}{\textbf{Se valora más\ldots}} &
  \textcolor{white}{\textbf{\ldots que}} \\
  \midrule
  Individuos e interacciones & Procesos y herramientas \\
  Software funcionando       & Documentación exhaustiva \\
  Colaboración con el cliente & Negociación contractual \\
  Respuesta ante el cambio   & Seguir un plan fijo \\
  \bottomrule
\end{tabular}

\vspace{0.4cm}
Las Historias de Usuario nacen directamente de esta filosofía: son la forma
ágil de capturar requisitos de manera simple, colaborativa y orientada
al valor para el usuario \parencite{beck2001}.

\subsection{Scrum como marco de trabajo}

\textbf{Scrum} es el marco ágil más utilizado en la industria
\parencite{schwaber2020}. Define los siguientes artefactos relacionados
con las Historias de Usuario:

\vspace{0.3cm}
\begin{tabular}{>{\bfseries}p{3.5cm} p{9cm}}
  \toprule
  \rowcolor{espeazul}
  \textcolor{white}{Artefacto} & \textcolor{white}{Descripción} \\
  \midrule
  Product Backlog & Lista priorizada de todas las Historias de Usuario del proyecto \\
  Sprint Backlog  & Subconjunto de historias seleccionadas para el sprint actual \\
  Incremento      & El producto funcional y verificable entregado al final de cada sprint \\
  \bottomrule
\end{tabular}

% =============================================================
\section{Historias de Usuario}

\subsection{Origen e historia}

Las Historias de Usuario fueron propuestas por \textbf{Kent Beck} como parte de
\textit{Extreme Programming (XP)} a finales de los años 90.
\textcite{jeffries2001} formalizó su estructura con el modelo 3C, y
posteriormente \textcite{cohn2004} las popularizó en su libro
\textit{User Stories Applied}, convirtiéndolas en el estándar de facto para
la captura de requisitos en entornos ágiles.

\subsection{¿Qué es una Historia de Usuario?}

Una \textbf{Historia de Usuario} es una descripción breve y simple de una
funcionalidad narrada desde la perspectiva del usuario final, cuyo propósito
es capturar \textbf{quién necesita qué y por qué}
\parencite{cohn2004}.

\textbf{Características fundamentales:}
\begin{itemize}[noitemsep, leftmargin=1.5cm]
  \item Se escriben en lenguaje natural, no técnico.
  \item Son breves y caben en una tarjeta de índice.
  \item Representan una unidad de valor para el usuario.
  \item Son el punto de partida de una conversación, no su fin.
  \item Deben poder completarse en un sprint (1 a 4 semanas).
\end{itemize}

\subsection{Formato estándar}

\begin{ejemplobox}
  \textit{Como \textbf{[ROL]}, quiero \textbf{[ACCIÓN]}, para \textbf{[BENEFICIO]}.}

  \vspace{0.3cm}
  \textbf{Ejemplo:} \textit{Como \textbf{estudiante}, quiero
  \textbf{subir un archivo PDF} como evidencia de una actividad, para
  \textbf{demostrar al profesor que completé la tarea asignada}.}
\end{ejemplobox}

\vspace{0.3cm}
\begin{tabular}{>{\bfseries}p{2.5cm} p{5cm} p{5cm}}
  \toprule
  \rowcolor{espeazul}
  \textcolor{white}{Elemento} & \textcolor{white}{Descripción} & \textcolor{white}{Pregunta que responde} \\
  \midrule
  Rol       & Quién usa la funcionalidad & ¿Para quién? \\
  Acción    & Qué quiere hacer           & ¿Qué necesita? \\
  Beneficio & Por qué lo necesita        & ¿Para qué sirve? \\
  \bottomrule
\end{tabular}

\subsection{Componentes: modelo 3C}

\textcite{jeffries2001} propone que una Historia de Usuario completa
se compone de tres partes:

\begin{itemize}[leftmargin=1.5cm]
  \item \textbf{Card (Tarjeta):} La descripción corta en el formato estándar.
        Es un recordatorio de la conversación, no una especificación completa.
  \item \textbf{Conversation (Conversación):} El diálogo continuo entre el
        equipo y el cliente para aclarar detalles y supuestos.
  \item \textbf{Confirmation (Confirmación):} Los \textbf{criterios de
        aceptación} que definen cuándo la historia está terminada y verificada.
\end{itemize}

\subsubsection{Criterios de Aceptación}

Los criterios de aceptación son condiciones medibles y verificables que debe
cumplir la funcionalidad para considerarse completada \parencite{wiegers2013}.
Deben ser acordados entre el Product Owner y el equipo antes de iniciar el
desarrollo.

\begin{ejemplobox}
  \textbf{HU-11:} Como estudiante, quiero subir un PDF como evidencia de una
  actividad, para demostrar que completé la tarea.

  \vspace{0.3cm}\textbf{Criterios de aceptación:}
  \begin{itemize}[noitemsep]
    \item Solo se aceptan archivos en formato PDF.
    \item El archivo no debe superar los 10 MB.
    \item El botón de subida solo aparece si el plazo no ha vencido.
    \item Al subir correctamente, el indicador cambia a ``Entregada''.
    \item Si supera el límite, el sistema muestra un mensaje de error descriptivo.
    \item Si el plazo venció, el botón aparece deshabilitado con ``Plazo cerrado''.
  \end{itemize}
\end{ejemplobox}

\subsection{Criterios INVEST}

Para que una Historia de Usuario sea de calidad, debe cumplir los criterios
\textbf{INVEST} \parencite{cohn2004}:

\vspace{0.3cm}
\begin{tabular}{>{\bfseries\centering\arraybackslash}p{1.1cm} >{\bfseries}p{3cm} p{5cm} p{3.5cm}}
  \toprule
  \rowcolor{espeazul}
  \textcolor{white}{Letra} & \textcolor{white}{Criterio} &
  \textcolor{white}{Significado} & \textcolor{white}{Señal de problema} \\
  \midrule
  I & Independent & No depende de otras historias & ``Necesita primero la HU-X'' \\
  N & Negotiable  & Los detalles son ajustables   & ``Debe ser exactamente así'' \\
  V & Valuable    & Genera valor al usuario       & ``El usuario no lo nota'' \\
  E & Estimable   & Se puede estimar su esfuerzo  & ``No sabemos cuánto toma'' \\
  S & Small       & Cabe en un sprint             & ``Tomaría meses'' \\
  T & Testable    & Se puede verificar            & ``No hay forma de saber'' \\
  \bottomrule
\end{tabular}

\subsection{Jerarquía: Épicos, Historias y Tareas}

Las historias se organizan de forma jerárquica dentro del proyecto
\parencite{schwaber2020}:

\begin{center}
  \textbf{Tema} $\supset$ \textbf{Épico} $\supset$ \textbf{Historia} $\supset$ \textbf{Tarea}
\end{center}

\vspace{0.3cm}
\begin{tabular}{>{\bfseries}p{2cm} p{5cm} p{5.5cm}}
  \toprule
  \rowcolor{espeazul}
  \textcolor{white}{Nivel} & \textcolor{white}{Descripción} & \textcolor{white}{Ejemplo en Gestión de Notas} \\
  \midrule
  Tema    & Área de negocio amplia (varios meses)     & ``Gestión Académica'' \\
  Épico   & Funcionalidad grande (varios sprints)     & ``Calificaciones y Promedios'' \\
  Historia & Funcionalidad pequeña (un sprint)         & ``HU-13: Calcular promedio ponderado'' \\
  Tarea   & Trabajo técnico (horas)                  & ``Crear endpoint POST /calificaciones'' \\
  \bottomrule
\end{tabular}

\subsection{Story Points y Planning Poker}

Los \textbf{Story Points} son una unidad relativa de esfuerzo y complejidad.
No miden horas sino la dificultad percibida en relación a otras historias.
Se utiliza la \textbf{escala Fibonacci} porque refleja la incertidumbre
creciente en tareas de mayor tamaño \parencite{cohn2004}.

\begin{formulabox}
  \textbf{Escala Fibonacci:} 1 → 2 → 3 → 5 → 8 → 13 → 21 → ?

  \vspace{0.3cm}
  \begin{tabular}{cp{9cm}}
    \toprule
    \rowcolor{espeazul}
    \textcolor{white}{\textbf{Puntos}} & \textcolor{white}{\textbf{Interpretación}} \\
    \midrule
    1  & Trivial — pocas horas de trabajo \\
    3  & Moderado — aproximadamente 1 día \\
    5  & Complejo — 2 días de trabajo \\
    8  & Muy complejo — 3 a 4 días \\
    13 & Enorme — una semana completa \\
    \bottomrule
  \end{tabular}
\end{formulabox}

\subsection{Errores comunes al escribir Historias de Usuario}

\begin{tabular}{p{3cm} p{5cm} p{5cm}}
  \toprule
  \rowcolor{espeazul}
  \textcolor{white}{\textbf{Error}} & \textcolor{white}{\textbf{Ejemplo incorrecto}} & \textcolor{white}{\textbf{Corrección}} \\
  \midrule
  Incluir solución técnica & ``Como usuario, quiero un endpoint REST\ldots'' & ``Como usuario, quiero iniciar sesión\ldots'' \\
  Historia muy grande      & ``Quiero gestionar todo el sistema''            & Dividir en historias específicas \\
  Sin beneficio claro      & ``Quiero un botón azul''                        & ``\ldots para identificarlo rápidamente'' \\
  Criterios vagos          & ``El sistema debe ser rápido''                  & ``Debe cargar en menos de 2 segundos'' \\
  Rol genérico             & ``Como usuario\ldots'' (siempre igual)          & Diferenciar: admin, profesor, estudiante \\
  \bottomrule
\end{tabular}

% =============================================================
\section{Puntos de Función (Function Points)}

\subsection{Origen e historia}

Los \textbf{Puntos de Función} fueron introducidos por
\textcite{albrecht1979} en IBM como respuesta a la necesidad de medir el
tamaño del software de forma independiente a la tecnología. Antes de los FP,
el único indicador de tamaño eran las \textbf{Líneas de Código (LOC)}, que
variaban enormemente entre lenguajes de programación.

\vspace{0.3cm}
\begin{tabular}{>{\centering\arraybackslash}p{1.5cm} p{11cm}}
  \toprule
  \rowcolor{espeazul}
  \textcolor{white}{\textbf{Año}} & \textcolor{white}{\textbf{Hito}} \\
  \midrule
  1979 & Albrecht introduce los Function Points en IBM \\
  1984 & Se funda el IFPUG (International Function Point Users Group) \\
  1994 & Se publica la primera versión del estándar ISO \\
  1999 & Aparece COSMIC como método alternativo para tiempo real \\
  2009 & ISO/IEC 20926 reconoce el método IFPUG oficialmente \parencite{iso2009} \\
  2010 & IFPUG publica la versión 4.3.1 del manual de conteo \parencite{ifpug2010} \\
  \bottomrule
\end{tabular}

\subsection{¿Qué son los Puntos de Función?}

Los \textbf{Puntos de Función (FP)} miden el \textbf{tamaño funcional} de un
sistema desde la perspectiva del usuario, independientemente del lenguaje,
la arquitectura o la plataforma \parencite{ifpug2010}.

\textbf{Características principales:}
\begin{itemize}[noitemsep, leftmargin=1.5cm]
  \item \textbf{Independiente de la tecnología:} los FP no cambian si el sistema
        se desarrolla en Java, Python o Node.js.
  \item \textbf{Orientado al usuario:} mide lo que el usuario percibe, no el
        código que escribe el programador.
  \item \textbf{Estable:} no cambia cuando se refactoriza el código interno.
  \item \textbf{Comparable:} permite comparar proyectos de distintas empresas y épocas.
  \item \textbf{Estandarizado:} reconocido por la norma ISO/IEC 20926 \parencite{iso2009}.
\end{itemize}

\subsection{Métodos de medición de tamaño funcional}

Existen varios métodos reconocidos internacionalmente:

\vspace{0.3cm}
\begin{tabular}{>{\bfseries}p{2.5cm} p{2cm} p{4.5cm} p{3cm}}
  \toprule
  \rowcolor{espeazul}
  \textcolor{white}{Método} & \textcolor{white}{Org.} &
  \textcolor{white}{Mejor para} & \textcolor{white}{Complejidad} \\
  \midrule
  IFPUG FPA & IFPUG & Sistemas de información web y empresariales & Media \\
  COSMIC FFP & COSMIC & Sistemas en tiempo real y embebidos & Alta \\
  NESMA & NESMA & Estimación rápida en fases tempranas & Baja \\
  FiSMA & Finnish SW & Sistemas pequeños y medianos & Media \\
  \bottomrule
\end{tabular}

\begin{notabox}
  Para sistemas web como \textbf{Gestión de Notas} (Next.js + NestJS + PostgreSQL),
  el método \textbf{IFPUG} es el más apropiado por su orientación a sistemas
  de información con interfaces de usuario, bases de datos y procesos de negocio
  \parencite{ifpug2010}.
\end{notabox}

\subsection{Tipos de componentes funcionales IFPUG}

\subsubsection{Funciones de datos — qué almacena el sistema}

\begin{tabular}{>{\bfseries}p{2.5cm} p{1cm} p{5cm} p{4cm}}
  \toprule
  \rowcolor{espeazul}
  \textcolor{white}{Tipo} & \textcolor{white}{Sigla} &
  \textcolor{white}{Descripción} & \textcolor{white}{Ejemplo en el proyecto} \\
  \midrule
  Arch. Lóg. Internos & ILF & Datos mantenidos por el propio sistema & Tablas: Usuarios, Notas, Evidencias \\
  Arch. Interfaz Ext. & EIF & Datos usados pero no mantenidos & APIs externas, BD de terceros \\
  \bottomrule
\end{tabular}

\subsubsection{Funciones de transacción — qué procesa el sistema}

\begin{tabular}{>{\bfseries}p{2.5cm} p{1cm} p{5cm} p{4cm}}
  \toprule
  \rowcolor{espeazul}
  \textcolor{white}{Tipo} & \textcolor{white}{Sigla} &
  \textcolor{white}{Descripción} & \textcolor{white}{Ejemplo en el proyecto} \\
  \midrule
  Entradas Externas  & EI & Datos que ingresan y modifican ILFs & Login, subida PDF, registrar nota \\
  Salidas Externas   & EO & Datos con lógica de procesamiento   & Promedio ponderado, PDF de notas \\
  Consultas Externas & EQ & Datos recuperados sin transformar   & Listado de alumnos, vista de notas \\
  \bottomrule
\end{tabular}

\subsection{Tabla de pesos por complejidad (IFPUG, 2010)}

\begin{table}[h!]
  \centering
  \caption{Pesos estándar IFPUG por tipo de componente y nivel de complejidad}
  \vspace{0.3cm}
  \begin{tabular}{>{\bfseries}p{5cm} >{\centering\arraybackslash}p{2cm} >{\centering\arraybackslash}p{2cm} >{\centering\arraybackslash}p{2cm}}
    \toprule
    \rowcolor{espeazul}
    \textcolor{white}{Tipo} & \textcolor{white}{Baja} & \textcolor{white}{Media} & \textcolor{white}{Alta} \\
    \midrule
    EI — Entradas Externas    & 3 & 4 & 6  \\
    EO — Salidas Externas     & 4 & 5 & 7  \\
    EQ — Consultas Externas   & 3 & 4 & 6  \\
    ILF — Arch. Lóg. Internos & 7 & 10 & 15 \\
    EIF — Arch. Interfaz Ext. & 5 & 7 & 10 \\
    \bottomrule
  \end{tabular}
\end{table}

\subsection{Proceso de cálculo paso a paso}

\subsubsection{Paso 1 — Definir el alcance}
Determinar qué módulos o funcionalidades se incluyen en la medición.

\subsubsection{Paso 2 — Identificar y clasificar los componentes}
Para cada historia o módulo, listar todos los EI, EO, EQ e ILF con su
complejidad \parencite{ifpug2010}.

\begin{ejemplobox}
  \textbf{HU-13:} \textit{``Como sistema, quiero calcular el promedio ponderado
  de cada parcial (TAREA 20\%, PRUEBA 20\%, PROYECTO 25\%, EXAMEN 35\%).''

  \vspace{0.3cm}
  \begin{tabular}{lcp{7.5cm}}
    \toprule
    Componente & Cantidad & Justificación \\
    \midrule
    EI  & 2 & Entrada de calificaciones + configuración de pesos \\
    EO  & 3 & Promedio del parcial + promedio final + estado del estudiante \\
    EQ  & 2 & Consulta de actividades + consulta de notas existentes \\
    ILF & 3 & Tablas: Calificaciones, Promedios, Actividades \\
    \bottomrule
  \end{tabular}}
\end{ejemplobox}

\subsubsection{Paso 3 — Calcular los PF No Ajustados (UFP)}

\begin{formulabox}
  \[
    UFP = \sum(EI \times w) + \sum(EO \times w) + \sum(EQ \times w) + \sum(ILF \times w)
  \]

  \textbf{Para HU-13 con complejidad Alta:}
  \[
    UFP = (2\times6) + (3\times7) + (2\times6) + (3\times15) = 12+21+12+45 = 90
  \]
  \textit{(Ajustado a 66 al refinar complejidades mixtas por componente.)}
\end{formulabox}

\subsubsection{Paso 4 — Factor de Ajuste de Valor (VAF)}

El VAF permite ajustar el resultado según 14 \textbf{Características Generales
del Sistema (GSC)}, cada una valorada de 0 (sin influencia) a 5 (esencial)
\parencite{ifpug2010}:

\vspace{0.3cm}
\begin{tabular}{>{\centering\arraybackslash}p{0.7cm} p{5.5cm} >{\centering\arraybackslash}p{0.7cm} p{5.5cm}}
  \toprule
  \rowcolor{espeazul}
  \textcolor{white}{\#} & \textcolor{white}{Característica} &
  \textcolor{white}{\#} & \textcolor{white}{Característica} \\
  \midrule
  1 & Comunicación de datos     & 8  & Actualización en línea \\
  2 & Procesamiento distribuido & 9  & Procesamiento complejo \\
  3 & Rendimiento               & 10 & Reusabilidad \\
  4 & Configuración utilizada   & 11 & Facilidad de instalación \\
  5 & Tasa de transacciones     & 12 & Facilidad de operación \\
  6 & Entrada de datos en línea & 13 & Múltiples sitios \\
  7 & Eficiencia del usuario    & 14 & Facilidad de cambios \\
  \bottomrule
\end{tabular}

\begin{formulabox}
  \[
    VAF = 0{,}65 + (0{,}01 \times \sum GSC) \qquad \text{Rango: } [0{,}65 \; ; \; 1{,}35]
  \]
  \[
    PF_{\text{Ajustados}} = UFP \times VAF
  \]
\end{formulabox}

\subsubsection{Paso 5 — Total del sistema y conversión a esfuerzo}

\begin{formulabox}
  \[
    PF_{\text{Total}} = \sum_{i=1}^{n} PF_i = 501 \text{ PF (Sistema Gestión de Notas)}
  \]
  \[
    \text{Esfuerzo} = 501 \times 6\,\frac{\text{h}}{\text{PF}} = 3.006\text{ horas}
    \qquad
    \text{Costo} = 3.006 \times \$20 = \$60.120\text{ USD}
  \]
\end{formulabox}

\subsection{Ventajas y limitaciones}

\begin{tabular}{p{6.5cm} p{6cm}}
  \toprule
  \rowcolor{espeazul}
  \textcolor{white}{\textbf{Ventajas}} & \textcolor{white}{\textbf{Limitaciones}} \\
  \midrule
  Independiente del lenguaje y la tecnología & Subjetividad en la complejidad \\
  Orientado al usuario final & Requiere entrenamiento certificado \\
  Comparable entre proyectos distintos & No mide calidad del código \\
  Estándar ISO/IEC 20926 \parencite{iso2009} & No apto para sistemas embebidos \\
  Base para COCOMO II y contratos & Proceso largo en sistemas grandes \\
  \bottomrule
\end{tabular}

% =============================================================
\section{Relación entre Historias de Usuario y Puntos de Función}

Las Historias de Usuario y los Puntos de Función son técnicas complementarias:
las primeras capturan \textbf{qué} debe hacer el sistema, mientras que los
segundos \textbf{cuantifican el tamaño} de lo que debe hacerse
\parencite{cohn2004, ifpug2010}.

\vspace{0.3cm}
\begin{tabular}{p{6.5cm} p{6.5cm}}
  \toprule
  \rowcolor{espeazul}
  \textcolor{white}{\textbf{Historia de Usuario provee\ldots}} &
  \textcolor{white}{\textbf{Puntos de Función mide\ldots}} \\
  \midrule
  Datos que ingresan (formularios, uploads) & $\rightarrow$ Entradas Externas (EI) \\
  Resultados calculados o reportes          & $\rightarrow$ Salidas Externas (EO) \\
  Consultas del usuario al sistema          & $\rightarrow$ Consultas Externas (EQ) \\
  Entidades en la base de datos             & $\rightarrow$ Arch. Lógicos Internos (ILF) \\
  Complejidad de criterios de aceptación    & $\rightarrow$ Nivel Baja / Media / Alta \\
  \bottomrule
\end{tabular}

\vspace{0.5cm}
\textbf{Comparación entre Story Points y Puntos de Función:}

\vspace{0.3cm}
\begin{tabular}{p{4cm} p{4.5cm} p{4.5cm}}
  \toprule
  \rowcolor{espeazul}
  \textcolor{white}{\textbf{Aspecto}} &
  \textcolor{white}{\textbf{Story Points}} &
  \textcolor{white}{\textbf{Puntos de Función}} \\
  \midrule
  Creado por        & Equipo (consenso)          & Analista certificado IFPUG \\
  Base              & Esfuerzo relativo          & Tamaño funcional objetivo \\
  Independiente del equipo & No               & Sí \\
  Uso principal     & Planificación de sprints   & Contratos y estimación de costos \\
  Estandarización   & No (escala propia)         & Sí (ISO/IEC 20926) \\
  Velocidad         & Rápida (Planning Poker)    & Lenta (análisis detallado) \\
  \bottomrule
\end{tabular}

% =============================================================
\section{Aplicación Práctica: Caso de Estudio}

\subsection{Sistema Gestión de Notas — Plataforma para Instituciones Educativas}

\begin{tabular}{ll}
  \textbf{Tipo de sistema:} & Plataforma genérica para instituciones educativas \\
  \textbf{Stack tecnológico:} & Next.js · NestJS · PostgreSQL · Prisma ORM \\
  \textbf{Período de desarrollo:} & 6 de abril al 5 de mayo de 2026 (29 días) \\
  \textbf{Sprints:} & 6 sprints de 5 días · 28 historias · 162 Story Points \\
  \textbf{Total PF:} & 501 Puntos de Función (16 historias analizadas) \\
\end{tabular}

\subsection{Resumen de PF por módulo}

\begin{table}[h!]
  \centering
  \caption{Distribución de Puntos de Función por módulo del sistema}
  \vspace{0.3cm}
  \begin{tabular}{p{8cm} >{\centering\arraybackslash}p{2.5cm} >{\centering\arraybackslash}p{2cm}}
    \toprule
    \rowcolor{espeazul}
    \textcolor{white}{\textbf{Módulo}} &
    \textcolor{white}{\textbf{Historias}} &
    \textcolor{white}{\textbf{PF}} \\
    \midrule
    Autenticación y Control de Roles    & HU-01, HU-02            & 20  \\
    CRUD Administrativo                 & HU-03 a HU-06           & 111 \\
    Gestión Transaccional Docente       & HU-07 a HU-09           & 116 \\
    Gestión Transaccional Estudiante    & HU-10 a HU-12           & 80  \\
    Motor de Automatización             & HU-13 a HU-15           & 147 \\
    Mantenimiento (Cron Job)            & HU-16                   & 27  \\
    \midrule
    \rowcolor{espeazul}
    \textcolor{white}{\textbf{TOTAL DEL SISTEMA}} &
    \textcolor{white}{\textbf{16 historias}} &
    \textcolor{white}{\textbf{501}} \\
    \bottomrule
  \end{tabular}
\end{table}

\subsection{Hallazgo principal}

El análisis IFPUG confirmó que el módulo de \textbf{Calificaciones y Promedios}
(HU-13 y HU-14) es el de mayor complejidad con \textbf{114 PF}, coincidiendo
con el sprint de mayor carga (Sprint 5, 29 story points). Esta coincidencia
entre la estimación ágil y la métrica formal valida la coherencia del
plan de desarrollo adoptado.

% =============================================================
\section{Conclusiones}

\begin{enumerate}[leftmargin=1.8cm, itemsep=4pt]
  \item Las \textbf{Historias de Usuario} son una herramienta fundamental en el
        desarrollo ágil porque centran el diseño en el valor para el usuario,
        facilitan la comunicación y permiten planificar el trabajo en sprints de
        forma incremental \parencite{cohn2004}.

  \item Los \textbf{Puntos de Función} complementan a las historias aportando
        una medida objetiva, estandarizada e independiente de la tecnología,
        que permite derivar estimaciones de esfuerzo y costo comparables
        entre proyectos \parencite{ifpug2010}.

  \item La \textbf{relación entre ambas técnicas} es directa: los elementos de
        una Historia de Usuario se corresponden uno a uno con los componentes
        funcionales IFPUG, facilitando la transición de la captura de requisitos
        a la estimación formal \parencite{sommerville2016}.

  \item En el proyecto \textbf{Gestión de Notas}, la combinación de Story Points
        (planificación interna) y Puntos de Función (estimación formal) permitió
        identificar el módulo más complejo del sistema y validar la distribución
        del esfuerzo a lo largo de los 6 sprints.

  \item Se recomienda el uso conjunto de ambas técnicas: Story Points para la
        planificación ágil interna y Puntos de Función para la comunicación
        formal con clientes y la gestión de contratos de software.
\end{enumerate}

% =============================================================
%  REFERENCIAS — generadas automáticamente por BibLaTeX/APA
% =============================================================
\newpage
\printbibliography[title={Referencias}]

\vfill
\begin{center}
  \small\color{grisoscuro}
  Construcción y Evolución del Software \quad|\quad
  ESPE · NRC 30744 \quad|\quad
  Luis Cueva · 2026
\end{center}

\end{document}
